\documentclass[10pt,a4paper]{amsart}
\usepackage[margin=19mm]{geometry}
\usepackage{amsmath,amssymb,mathtools}
\usepackage[scr=rsfs,cal=euler]{mathalfa}
\usepackage{xfrac}
\usepackage[unicode,hidelinks,bookmarks=false]{hyperref}
\newcommand{\ie}{i.e.\ }
\newcommand{\cf}{cf.\ }
\newcommand{\resp}{respectively\ }
\newcommand{\Subsection}{\subsection}
\providecommand{\nobreakdash}{\nobreak\hbox{-}}
\setlength{\emergencystretch}{2em}
\multlinegap0pt
\usepackage{bm}
\usepackage{bbm}
\usepackage{dsfont}
\usepackage{xspace}
\usepackage{cancel}
\usepackage{mathtools}
\usepackage{amsthm}
\makeatletter
\@ifundefined{teorema}{\newtheorem{teorema}{Teorema}}{}
\@ifundefined{lema}{\newtheorem{lema}[teorema]{Lemma}}{}
\makeatother
\DeclareMathOperator{\Ker}{Ker}
\title[Addendum/Corrigendum to "On solubility of skew left braces a]{Addendum/Corrigendum to "On solubility of skew left braces and solutions of the Yang-Baxter equation"}
\author{A. Ballester-Bolinches, R. Esteban-Romero, P. Jim\'enez-Seral, V. P\'erez-Calabuig}
\date{Source: arXiv:2604.20894v1 (CC BY 4.0)}
\begin{document}
\maketitle

\noindent\emph{Source and licence.} Addendum/Corrigendum to "On solubility of skew left braces and solutions of the Yang-Baxter equation". Version arXiv:2604.20894v1 (CC BY 4.0), distributed under the Creative Commons Attribution 4.0 International licence, as recorded in the arXiv licence metadata for this exact version and shown on the abstract page https://arxiv.org/abs/2604.20894v1. Full rights notice: licenses/arxiv-2604.20894-CC-BY-4.0.txt.\par\smallskip

\noindent\textit{Editorial scope.} Counted unit: the Lema whose source label is \texttt{lema:kerbarf} in the article (source0.tex lines 300--304, source label \texttt{lema:kerbarf}), delivered together with the article's own complete original proof of it (source0.tex lines 306--340, one \texttt{proof} environment of 35 lines). No mathematics is written, repaired or re-derived: statement and proof are the article's own text line for line. Required context retained uncounted: eq:x(ker f)y->Kerf (lines 292--295). Every definition, display and label of those lines is retained unchanged. Omitted from this package and not redistributed: 1--291, 296--299, 305--305, 341--910. Nothing inside a retained excerpt is omitted or altered: every retained body is delivered exactly as the article's lines are written, so no omission receipt of any kind accompanies this package. Citations: the retained excerpts carry no citation command at all, so no bibliography entry of the article is transcribed and no reference is redistributed. Reference adaptation: none was needed and none was made; every reference inside the delivered unit resolves inside it, so no unresolved reference or citation is produced. Printed slips kept literal: no printed word, symbol, hypothesis or number of the counted statement or of its proof was altered. Layout and class: the article's own class is \texttt{article} with packages including fontenc, babel, fancyhdr, inputenc, amsmath, amsthm, amsfonts, amssymb, enumerate, graphicx, indentfirst, url, xy; the wrapper is the lane's pinned \texttt{amsart} editorial wrapper at 10pt on A4, which restates the theorem-like environments the retained text needs and adds the article's own macro definitions that the retained text needs (\texttt{\textbackslash Ker}), with extra wrapper packages (bm, bbm, dsfont, xspace, cancel, mathtools). The delivered numbering is the unit's own. Assets: no retained span contains a figure environment, a graphics-inclusion command, a PDF-page inclusion, a table or a TikZ picture; no image file is copied into this package, the package declares no asset dependency and no image file is redistributed with it. Licence: version arXiv:2604.20894v1 of this article is distributed under the Creative Commons Attribution 4.0 International licence, as recorded in the arXiv licence metadata for this exact version and shown on the abstract page https://arxiv.org/abs/2604.20894v1. A full rights notice is delivered at licenses/arxiv-2604.20894-CC-BY-4.0.txt. Characters: every retained excerpt is pure ASCII, so the delivered engine, XeLaTeX with its default Latin Modern text font, reports no missing character.
\begin{equation}
\label{eq:x(ker f)y->Kerf}
f(x) = f(y) \Rightarrow \bar{f}(x^{-1}y) = \bar{f}(x)^{-1}\bar f(y) = f(x)^{-1}f(y) = 0 \in G_Y
\end{equation}

\begin{lema}
\label{lema:kerbarf}
Let $f \colon X \rightarrow Y$ be an epimorphism of solutions $(X,r)$ and $(Y,s)$, with induced homomorphism $\bar f\colon G_X \rightarrow G_Y$. Then, $\Ker \bar f$ is the ideal generated 
\[ K:= \{x^{-1}y \mid \text{$x,y \in X$ with $x (\ker f) y$}\}.\]
\end{lema}

\begin{proof}
Write $I_K$ the ideal in $G_X$ generated by~$K$. Clearly, \eqref{eq:x(ker f)y->Kerf} yields one inclusion. For the other inclusion, take 
\[ R_{\mathsf{X}} := \{ \mathsf{u}\mathsf{v}\rho_{\mathsf{v}}^{\mathsf{r}}(\mathsf{u})^{-1}\lambda_{\mathsf{u}}^{\mathsf{r}}(\mathsf v)^{-1} \mid \mathsf{u,v} \in \mathsf{X}\},\] for each $\mathsf{X}\in \{X,Y\}$ with respectively $\mathsf{r}\in \{r,s\}$. Since $f$ is surjective and it makes $\lambda$ and $\rho$ maps of $r$ and $s$ compatible, it follows that for every $w\in R_Y$, there exists $u \in R_X$ such that $\bar f(u) = w$.

%observe that elements in $I_K$ are recursively generated by the following rules
%\begin{itemize}
%\item for every $w\in K$, $w \in I_K$;
%\item if $w \in I_K$, then $w^{-1}, -w, \lambda_{w'}(w), w'+w-w', w'w(w')^{-1}\in I_K$ for every $w'\in G_X$;
%\item if $w_1,w_2\in I_K$, then $w_1w_2^{-1}, w_1 - w_2 \in I_K$.
%\end{itemize}


Let $w_0 \in \Ker \bar f$ and write $w_0 = x_1^{\varepsilon_1}\cdots x_n^{\varepsilon_n}$ for some $x_i \in X$, $\varepsilon_i =\pm 1$, $1\leq i \leq n$. Then,
\[ 0 = \bar f(w_0) = f(x_1)^{\varepsilon_1}\cdots f(x_n)^{\varepsilon_n} \in G_Y.\]
Thus, there exist $v_1, \ldots, v_m \in G_Y$ and $w_1, \ldots, w_m \in R_Y$ such that the word
\[
w':= f(x_1)^{\varepsilon_1} \cdots f(x_n)^{\varepsilon_n} (v_1 w_1 v_1^{-1}) \cdots (v_m w_m v_m^{-1})
\]
reduces recursively to $0$
%\[
%\tilde{w}' = (\tilde{v}_1 \tilde{w}_1 \tilde{v}_1^{-1}) \cdots (\tilde{v}_k \tilde{w}_k \tilde{v}_k^{-1}),
%\]
%where $\tilde{v}_1, \ldots, \tilde{v}_k \in G_Y$ and $\tilde{w}_1, \ldots, \tilde{w}_k \in R_Y$, 
by successive cancellations of subwords of the form either $yy^{-1}$ or $y^{-1}y$, with $y \in Y$.

Since $f$ is surjective, for each $1 \leq i \leq m$, there exists a word $z_i \in G_X$ such that
\[
\mathrm{length}(z_i) = \mathrm{length}(v_i)\quad\text{and} \quad \bar f(z_i) = v_i,
\]
where each letter of $z_i$ maps under $\bar f$ to the corresponding letter of $v_i$. Moreover, for every $1\leq i \leq m$, there exists $u_i \in R_X$ such that $f(u_i)=w_i$.

Thus, $w_0':= w_0(z_1u_1z_1^{-1})\cdots (z_mu_mz_m^{-1}) = w_0 \in G_X$ and $\bar f(w_0') = w' = 0 \in G_Y$, where $\mathrm{length}(w_0') = \mathrm{length}(w')$, and each letter of $w_0'$ maps under $\bar f$ to the corresponding letter of $w'$. Therefore, in $G_X/I_K$, the word $w_0'I_K$ recursively reduces to $0I_K$ following the same reduction procedure than $w'$: there is a one-to-one correspondence between each cancellation of the form $y^{-1}y \in G_Y$ with $x^{-1}z \in X$, where $f(x) = f(z) = y$, and similarly for $yy^{-1}$.

Hence, $w_0 = w_0'\in G_X$ and $w_0' \in I_K$, so that the other inclusion also holds.
\end{proof}

\end{document}
